\documentclass[11pt,a4paper]{article}
\usepackage[utf8]{inputenc}
\usepackage[spanish]{babel}
\usepackage{amsmath, amssymb}
\usepackage{geometry}
\geometry{top=2.5cm, bottom=2.5cm, left=2.5cm, right=2.5cm}

\begin{document}

\section*{Exámenes de práctica de los capítulos 1, 2 y 3}

\subsection*{Examen 1}

\begin{itemize}
    \item[\textbf{1.}] De cierta roca uniforme se cortan dos esferas. Una tiene $4.50 \text{ cm}$ de radio. La masa de la segunda esfera es cinco veces mayor. Encuentre el radio de la segunda esfera.
    \vspace*{6cm}

    \item[\textbf{2.}] Un galón de pintura (volumen $= 3.78 \times 10^{-3} \text{ m}^3$) cubre un área de $25.0 \text{ m}^2$. ¿Cuál es el grosor de la pintura fresca sobre la pared?
    \vspace*{6cm}

    \item[\textbf{3.}] A partir del conjunto de ecuaciones
    \begin{align*}
        p &= 3q \\
        pr &= qs \\
        \frac{1}{2}pr^2 + \frac{1}{2}qs^2 &= \frac{1}{2}qt^2
    \end{align*}
    que contienen las incógnitas $p, q, r, s$ y $t$, encuentre el valor de la proporción de $t$ con $r$.
    \vspace*{7cm}

    \item[\textbf{4.}] La rapidez de un impulso nervioso en el cuerpo humano es de aproximadamente $100 \text{ m/s}$. Si su dedo del pie tropieza accidentalmente en la oscuridad, estime el tiempo que tarda el impulso nervioso en viajar a su cerebro.
    \vspace*{5cm}
\end{itemize}

\newpage

\begin{itemize}
    \item[\textbf{5.}] Un objeto que se mueve con aceleración uniforme tiene una velocidad de $12.0 \text{ cm/s}$ en la dirección $x$ positiva cuando su coordenada $x$ es $3.00 \text{ cm}$. Si su coordenada $x$, $2.00 \text{ s}$ después es $-5.00 \text{ cm}$, ¿cuál es su aceleración?
    \vspace*{6cm}

    \item[\textbf{6.}] Lisa llega de prisa a un andén del metro para encontrar que su tren está a punto de partir. Se detiene y mira los carros que pasan. Cada carro tiene una longitud de $8.60 \text{ m}$. El primero se mueve pasándola en $1.50 \text{ s}$ y el segundo en $1.10 \text{ s}$. Encuentre la aceleración constante del tren.
    \vspace*{6cm}

    \item[\textbf{7.}] Su perro está corriendo alrededor de la hierba en el patio trasero. Él experimenta desplazamientos sucesivos $3.50 \text{ m}$ sur, $8.20 \text{ m}$ al noreste y $15.0 \text{ m}$ al oeste. ¿Cuál es el desplazamiento resultante?
    \vspace*{6cm}

    \item[\textbf{8.}] Dos vectores $\vec{A}$ y $\vec{B}$ tienen magnitudes exactamente iguales. Para que la magnitud de $\vec{A} + \vec{B}$ sea $100$ veces más grande que la magnitud de $\vec{A} - \vec{B}$, ¿cuál debe ser el ángulo entre ellos?
    \vspace*{6cm}
\end{itemize}

\newpage

\subsection*{Examen 2}

\begin{itemize}
    \item[\textbf{1.}] ¿Cuáles de las siguientes ecuaciones son dimensionalmente correctas? (a) $v_f = v_i + ax$ o (b) $y = (2 \text{ m}) \cos(kx)$, donde $k = 2 \text{ m}^{-1}$.
    \vspace*{6cm}

    \item[\textbf{2.}] El \textit{año tropical}, el intervalo desde un equinoccio de primavera hasta el siguiente equinoccio de primavera, es la base para el calendario. Contiene $365.242\ 199$ días. Encuentre el número de segundos en un año tropical.
    \vspace*{6cm}

    \item[\textbf{3.}] La distancia del Sol a la estrella más cercana es casi de $4 \times 10^{16} \text{ m}$. La Vía Láctea (véase la siguiente figura) es en términos aproximados un disco de $\sim 10^{21} \text{ m}$ de diámetro y $\sim 10^{19} \text{ m}$ de grosor. Encuentre el orden de magnitud del número de estrellas en la Vía Láctea. Considere representativa la distancia entre el Sol y el vecino más cercano.
    \vspace*{6cm}

    \item[\textbf{4.}] Una persona que hace un viaje conduce con una rapidez constante de $89.5 \text{ km/h}$, excepto por una parada de descanso de $22.0 \text{ min}$. Si la rapidez promedio de la persona es de $77.8 \text{ km/h}$, (a) ¿cuánto tiempo invierte la persona en el viaje y (b) qué tan lejos llegará?
    \vspace*{6cm}
\end{itemize}

\newpage

\begin{itemize}
    \item[\textbf{5.}] La altura de un helicóptero sobre el suelo está dada por $h = 3.00t^3$, donde $h$ está en metros y $t$ en segundos. Después de $2.00 \text{ s}$, el helicóptero libera una pequeña valija de correo. ¿Cuánto tiempo después de su liberación, la valija llega al suelo?
    \vspace*{6cm}

    \item[\textbf{6.}] En una carrera de $100 \text{ m}$ de mujeres, a Laura, acelerando de manera uniforme, le toma $2.00 \text{ s}$ y a Healan $3.00 \text{ s}$ alcanzar sus rapideces máximas, que mantienen cada una durante el resto de la carrera. Cruzan la línea de meta al mismo tiempo, estableciendo ambas un récord mundial de $10.4 \text{ s}$. (a) ¿Cuál es la aceleración de cada atleta? (b) ¿Cuáles son sus respectivas rapideces máximas? (c) ¿Qué atleta está por delante en la marca de $6.00 \text{ s}$, y por cuánto? (d) ¿Cuál es la distancia máxima en que Healan está detrás de Laura, y en qué tiempo eso ocurre?
    \vspace*{8cm}

    \item[\textbf{7.}] Considere los tres vectores desplazamiento $\vec{A} = (3\hat{i} - 3\hat{j}) \text{ m}$, $\vec{B} = (\hat{i} - 4\hat{j}) \text{ m}$ y $\vec{C} = (-2\hat{i} + 5\hat{j}) \text{ m}$. Use el método de componentes para determinar (a) la magnitud y dirección del vector $\vec{D} = \vec{A} + \vec{B} + \vec{C}$ y (b) la magnitud y dirección de $\vec{E} = -\vec{A} - \vec{B} + \vec{C}$.
    \vspace*{6cm}

    \item[\textbf{8.}] Un vector está dado por $\vec{R} = 2\hat{i} + \hat{j} + 3\hat{k}$. Encuentre (a) las magnitudes de los componentes $x, y$ y $z$, (b) la magnitud de $\vec{R}$ y (c) los ángulos entre $\vec{R}$ y los ejes $x, y$ y $z$.
    \vspace*{6cm}
\end{itemize}

\end{document}
